\begin{center}
Problema B

Jack Lanterna e as trilhas

{\small Nome do arquivo: B.c, B.cpp, B.java, ou B.py}

{\large Timelimit: $1$}

{\tiny Por André da Cunha Ribeiro}

\end{center}					

{\footnotesize
Jack Lanterna é um caçador de monstros muito temido, que na noite de Halloween decide explorar uma cidade assombrada nos confins de uma floresta. A cidade possui casas mal-assombradas conectadas por trilhas, cada uma com um número de monstros que guardam o caminho.

Cada vez que Jack atravessa uma trilha e derrota os monstros que a guardam, ele destrói a trilha com sua magia, impedindo o retorno por aquele caminho. Jack consegue descansar em cada casa para recuperar suas forças antes de enfrentar mais monstros. Sua missão é derrotar todos os monstros e, ao final, retornar à casa inicial.

Você precisa determinar se Jack Lanterna consegue eliminar todos os monstros e retornar à casa inicial após derrotar todos os monstros.}

\vspace{0.3cm}
\textbf{Entrada}
A entrada contém vários casos de teste. A primeira linha de cada caso de teste contém dois inteiros $N$, $M$, indicando respectivamente o número de casas $(1 \leq N \leq 1000)$, e de trilhas $(1 \leq N \leq 1000)$. 

Cada uma das $M$ linhas seguintes contém três inteiros $u$, $v$ e $k$ que descrevem uma trilha entre as casas $u$ e $v$ e a quantidade de $k$ monstros nela.

\vspace{0.3cm}
\textbf{Saída}
Para cada caso de teste, o programa deve produzir uma linha na saída contendo $S$ se Jack Lanterna consegue eliminar todos os monstros e retornar à sua casa inicial, ou $N$ caso contrário.

{\tiny
\begin{table}[ht]
    \centering
    \begin{tabular}{|p{7cm}|p{7cm}|}
        \hline
        \begin{center}
            \vspace{-0.3cm}
            \textbf{Exemplos de Entrada}
            \vspace{-0.3cm}
        \end{center} 
         &  
        \begin{center}
            \vspace{-0.3cm}
            \textbf{Exemplos de Saída}
            \vspace{-0.3cm}
        \end{center} \\ \hline
         \texttt{5 6} & \texttt{S} \\  
         \texttt{1 2 1} & \texttt{} \\
         \texttt{1 3 2} & \texttt{} \\
         \texttt{1 4 3} & \texttt{} \\
         \texttt{1 5 4} & \texttt{} \\
         \texttt{2 3 2} & \texttt{} \\
         \texttt{4 5 1} & \texttt{} \\
         \hline
         \texttt{5 7} & \texttt{N} \\  
         \texttt{1 2 1} & \texttt{} \\
         \texttt{1 3 2} & \texttt{} \\
         \texttt{1 4 3} & \texttt{} \\
         \texttt{1 5 4} & \texttt{} \\
         \texttt{2 3 2} & \texttt{} \\
         \texttt{3 4 1} & \texttt{} \\
         \texttt{4 5 1} & \texttt{} \\
         \hline
         \texttt{6 6} & \texttt{N} \\  
         \texttt{1 2 1} & \texttt{} \\
         \texttt{1 3 2} & \texttt{} \\
         \texttt{2 3 3} & \texttt{} \\
         \texttt{4 5 1} & \texttt{} \\
         \texttt{5 6 2} & \texttt{} \\
         \texttt{6 4 3} & \texttt{} \\
        \hline
    \end{tabular}
    \caption{Exemplos de entradas e saídas}
    \label{tab:inputJ3}
\end{table}}

\newpage

\begin{center}
\noindent
\lstinputlisting{abobora_24/B/B4.cpp}
\end{center}